\documentclass{article}
\usepackage{fontspec}
\usepackage[margin=1in]{geometry}
\usepackage{amsmath, amssymb}
\usepackage{xcolor}
\usepackage[most]{tcolorbox}
\usepackage{tikz}
\usepackage{pgfplots}
\pgfplotsset{compat=1.18}
\usepackage{listings}
\usepackage{booktabs}
\usepackage{colortbl}
\usepackage{algorithm}
\usepackage{algpseudocode}
\usepackage{hyperref}
\title{Class / Lecture Notes}
\author{Generated by AI Meeting Pro}
\date{\today}
\begin{document}
\maketitle
## Core Concepts Explained

The lecture discusses the historical debate between wave and particle theories of light, focusing on the experiments of Thomas Young (double-slit experiment) and Albert Einstein (photoelectric effect). The experiments demonstrate that light exhibits both wave-like and particle-like behavior, leading to the development of quantum mechanics.

### Wave-Particle Duality

Light behaves as a wave in certain phenomena such as diffraction, refraction, and interference, while it behaves as particles (photons) in others like the photoelectric effect. This dual nature is known as wave-particle duality.

### Quantum Mechanics

Quantum mechanics is a fundamental theory that describes the behavior of matter and energy at the atomic and subatomic levels. It provides a framework for understanding phenomena such as the photoelectric effect, quantum entanglement, and particle-wave duality.

## Important Definitions & Formulas

1. Wave-Particle Duality: The property of particles (such as photons) to exhibit both wave-like and particle-like behavior.
2. Photon: A fundamental particle that is the quantum unit of light and other electromagnetic radiation.
3. Quantum Mechanics: A theory that describes the behavior of matter and energy at the atomic and subatomic levels, including wave-particle duality, quantum entanglement, and probabilistic nature of particles.
4. Probability Wave Function (Ψ): In quantum mechanics, a mathematical function that describes the probability distribution of a particle in space and time.
5. Quantum State: The complete set of properties and values of a quantum system at a given instant.
6. Superposition Principle: In quantum mechanics, the principle stating that any two (or more) quantum states can be added together to form a new state, and that each possible outcome is associated with an amplitude or probability.
7. Quantum Entanglement: A phenomenon in which pairs or groups of particles become linked and instantaneously affect each other's state no matter the distance between them.
8. Photoelectric Effect: The emission of electrons from a material when it is exposed to light, especially ultraviolet or X-rays. The kinetic energy of the emitted electrons depends on the frequency (color) of the incident light and a threshold frequency for electron emission.
9. Stopping Potential: In the photoelectric effect, the electric potential difference required to stop the motion of an emitted electron. It is equal to the kinetic energy of the electron.

## Study Guide / Key Takeaways

1. Light exhibits both wave-like and particle-like behavior (wave-particle duality).
2. Quantum mechanics is a fundamental theory that describes the behavior of matter and energy at the atomic and subatomic levels, including wave-particle duality, quantum entanglement, and probabilistic nature of particles.
3. The double-slit experiment by Thomas Young demonstrates the wave-like properties of light, while the photoelectric effect by Albert Einstein shows its particle-like behavior.
4. Quantum mechanics uses mathematical functions such as the probability wave function (Ψ) to describe the probability distribution of a particle in space and time.
5. The superposition principle states that any two quantum states can be added together to form a new state, and each possible outcome is associated with an amplitude or probability.
6. Quantum entanglement is a phenomenon in which pairs or groups of particles become linked and instantaneously affect each other's state no matter the distance between them.
7. The photoelectric effect involves the emission of electrons from a material when it is exposed to light, with the kinetic energy of the emitted electrons depending on the frequency (color) of the incident light and a threshold frequency for electron emission.
8. The stopping potential in the photoelectric effect is the electric potential difference required to stop the motion of an emitted electron, equal to the kinetic energy of the electron.

\vfill
\begin{center}
\footnotesize \textit{Generated by AI Scholar Hub on May 07, 2026 at 08:40 PM}
\end{center}

\end{document}